\chapter{Contratti e modello dei dati}
\label{cap:05-contratti-dati}

Il modulo \code{platform/common} contiene i tipi che attraversano ogni confine
del sistema: quelli che il control plane scambia con i chiamanti, con i runtime
delle funzioni e con gli SDK nei vari linguaggi. Sono meno di cinquecento righe
di codice, e la loro stabilit\`a \`e ci\`o che consente ai moduli e agli SDK di
evolvere indipendentemente. Questo capitolo li descrive e ne discute le scelte
non ovvie.

\section{\code{FunctionSpec}: la definizione di una funzione}

\code{FunctionSpec} \`e il tipo che il chiamante invia per registrare una
funzione e che il control plane conserva. \`E un \code{record} Java con sedici
componenti, raggruppabili in cinque famiglie.

\begin{center}
\small
\begin{adjustbox}{max width=\textwidth}
\begin{tabular}{@{}llp{6.2cm}@{}}
\toprule
\textbf{Famiglia} & \textbf{Campi} & \textbf{Ruolo} \\
\midrule
Identit\`a & \code{name}, \code{image} &
Obbligatori, non vuoti. Il nome \`e la chiave logica. \\
Esecuzione & \code{command}, \code{env}, \code{runtimeMode},
\code{runtimeCommand} &
Come avviare e come parlare con il processo funzione. \\
Risorse & \code{resources} &
Richieste e limiti di CPU e memoria, validati come coerenti fra loro. \\
Politiche & \code{timeoutMs}, \code{concurrency}, \code{queueSize},
\code{maxRetries} &
Sovrascrivono i default di piattaforma. \\
Collocazione & \code{executionMode}, \code{endpointUrl},
\code{imagePullSecrets}, \code{scalingConfig}, \code{offload} &
Dove la funzione vive e come viene governata. \\
\bottomrule
\end{tabular}
\end{adjustbox}
\end{center}

Tutti i campi delle ultime tre famiglie sono tipi \emph{boxed}
(\code{Integer}, non \code{int}). La scelta \`e deliberata e ricorre in tutto il
modello: il valore \code{null} significa ``non specificato'', ed \`e
distinguibile da qualunque valore esplicito. Un campo \code{int} con default
zero renderebbe indistinguibili ``l'utente non ha detto nulla'' e ``l'utente ha
chiesto zero'', e la risoluzione dei default (\cref{cap:06-nucleo}) non
potrebbe sapere quando intervenire.

\subsection{Validazione dichiarativa}

I vincoli sono espressi con annotazioni Jakarta Validation applicate direttamente
ai componenti del record:

\begin{lstlisting}[language=Java]
public record FunctionSpec(
        @NotBlank String name,
        @NotBlank String image,
        ...
        @Valid ResourceSpec resources,
        @Min(1) Integer timeoutMs,
        @Min(1) Integer concurrency,
        @Min(1) Integer queueSize,
        @Min(0) Integer maxRetries,
        ...
) { ... }
\end{lstlisting}

Il vincolo pi\`u interessante \`e quello che le annotazioni sui singoli campi non
possono esprimere, ossia la coerenza fra richiesta e limite di risorsa. \`E
codificato come metodo di validazione sul record che lo riguarda:

\begin{lstlisting}[language=Java, caption={Vincolo relazionale in
\code{ResourceSpec}.}]
@AssertTrue(message = "resource request must not exceed limit")
public boolean isRequestWithinLimit() {
    if (requests == null || limits == null) {
        return true;
    }
    return notGreater(requests.cpu(), limits.cpu())
            && notGreater(requests.memoryMiB(), limits.memoryMiB());
}
\end{lstlisting}

La conseguenza pratica \`e che una specifica incoerente viene respinta con
\code{400} al momento della registrazione, e non produce un deployment che
Kubernetes rifiuterebbe pi\`u tardi con un messaggio meno leggibile.

\subsection{Costruttori di compatibilit\`a}

\code{FunctionSpec} espone tre costruttori: quello canonico a sedici componenti
e due pi\`u corti che delegano al canonico passando \code{null} per i campi
aggiunti nelle evoluzioni successive (\code{offload} e, prima, \code{scalingConfig}).

\`E una tecnica esplicita di gestione della compatibilit\`a in un progetto che
usa i \code{record}: un record aggiunge un componente e il costruttore canonico
cambia firma, il che rompe ogni chiamante --- comprese, in questo progetto,
diverse centinaia di istanziazioni nei test. I costruttori di compatibilit\`a
assorbono la rottura. Il prezzo \`e che il tipo porta con s\'e la propria storia
di versioni, e il beneficio \`e che l'aggiunta di un campo resta un'operazione a
diff contenuto.

\section{\code{InvocationRequest} e \code{InvocationResponse}}

Il contratto di invocazione \`e volutamente povero.

\begin{lstlisting}[language=Java]
public record InvocationRequest(
        @NotNull(message = "Input payload is required") Object input,
        Map<String, String> metadata,
        Map<String, String> headers
) { ... }
\end{lstlisting}

Il campo \code{input} \`e di tipo \code{Object}: la piattaforma non impone uno
schema al payload e si limita a trasportarlo. \`E \code{@NotNull} ma pu\`o essere
qualunque struttura JSON. La distinzione fra \code{metadata} e \code{headers} \`e
di provenienza: \code{metadata} \`e ci\`o che il chiamante dichiara nel corpo,
\code{headers} \`e ci\`o che il control plane ha effettivamente osservato sulla
richiesta HTTP e ricostruisce prima di consegnarla
(\cref{cap:04-architettura-control-plane}).

La risposta trasporta l'esito e, opzionalmente, una \emph{busta} che la funzione
pu\`o usare per governare la risposta HTTP:

\begin{lstlisting}[language=Java]
public record InvocationResponse(
        String executionId, String status, Object output, ErrorInfo error,
        Integer statusCode, Map<String, String> headers, String encoding
) { ... }
\end{lstlisting}

I tre campi finali sono il meccanismo con cui una funzione pu\`o restituire, per
esempio, un \code{404} applicativo o un contenuto binario codificato in base64.
Il control plane li applica alla risposta soltanto se il codice di stato ricade
nell'intervallo valido, e in tal caso marca la risposta con
\code{X-NanoFaaS-Function-Status: true} cos\`i che il chiamante possa distinguere
un errore deciso dalla funzione da un errore deciso dalla piattaforma. \`E una
distinzione che le piattaforme FaaS spesso non rendono osservabile e la cui
assenza rende ambigua l'interpretazione di un tasso di errore.

\section{\code{ExecutionStatus} e il ciclo di vita}

\code{ExecutionStatus} \`e la proiezione esterna di un'esecuzione, restituita da
\code{GET /v1/executions/\{id\}}. Oltre all'esito porta due campi diagnostici
--- \code{coldStart} e \code{initDurationMs} --- che rendono osservabile per
singola invocazione se questa ha pagato un'inizializzazione e quanto. Sono la
sorgente della metrica \code{function\_cold\_start\_ms}
(\cref{cap:21-osservabilita}).

Lo stato \`e rappresentato come stringa e non come enumerazione, il che \`e una
scelta discutibile su cui vale la pena essere espliciti: agevola l'evoluzione del
vocabolario di stato senza rompere i client, ma sposta a runtime un controllo che
il compilatore avrebbe potuto fare. I valori usati sono \code{queued},
\code{running}, \code{succeeded}, \code{failed} e \code{timeout}.

\section{\code{InvocationResult}: il contratto verso il runtime}

Mentre \code{InvocationResponse} \`e ci\`o che il control plane restituisce al
chiamante, \code{InvocationResult} \`e ci\`o che il \emph{runtime della funzione}
restituisce al control plane --- sia come corpo della risposta HTTP sincrona,
sia come corpo della callback asincrona verso
\code{/v1/internal/executions/\{id\}:complete}.

\begin{lstlisting}[language=Java]
public record InvocationResult(
        boolean success, Object output, ErrorInfo error,
        Integer statusCode, Map<String, String> headers, String encoding
) {
    public static InvocationResult success(Object output) { ... }
    public static InvocationResult successWithEnvelope(Object output, Integer statusCode,
                                                       Map<String, String> headers, String encoding) { ... }
    public static InvocationResult error(String code, String message) { ... }
}
\end{lstlisting}

I metodi factory statici sono l'API che gli SDK usano
(\cref{cap:18-sdk}); l'esistenza di \code{successWithEnvelope} accanto a
\code{success} rende esplicito nel codice della funzione quando si sta
esercitando il controllo sulla risposta HTTP e quando no.

\section{Le enumerazioni di politica}

Cinque enumerazioni definiscono lo spazio delle politiche configurabili. Sono
elencate qui per completezza; ciascuna \`e trattata nel capitolo del modulo che
la implementa.

\begin{center}
\small
\begin{adjustbox}{max width=\textwidth}
\begin{tabular}{@{}lll@{}}
\toprule
\textbf{Enumerazione} & \textbf{Valori} & \textbf{Trattata in} \\
\midrule
\code{ExecutionMode} & \code{LOCAL}, \code{EXTERNAL}, \code{DEPLOYMENT} &
\cref{cap:07-modalita-esecuzione} \\
\code{RuntimeMode} & \code{HTTP}, \code{STDIO}, \code{FILE} &
\cref{cap:17-runtime-funzione} \\
\code{ScalingStrategy} & \code{HPA}, \code{INTERNAL}, \code{NONE} &
\cref{cap:11-autoscaler} \\
\code{ConcurrencyControlMode} & \code{FIXED}, \code{STATIC\_PER\_POD},
\code{ADAPTIVE\_PER\_POD}, \code{BUDGETED}, \code{SOJOURN} &
\cref{cap:12-concurrency-control} \\
\code{SyncQueueRejectReason} & \code{DEPTH}, \code{EST\_WAIT}, \code{TIMEOUT} &
\cref{cap:10-sync-queue} \\
\bottomrule
\end{tabular}
\end{adjustbox}
\end{center}

\section{Il blocco di scaling e la sua composizione}

\code{ScalingConfig} contiene la strategia di scaling, i limiti di repliche, una
lista di metriche e --- annidato --- il blocco di controllo della concorrenza:

\begin{lstlisting}[language=Java]
public record ScalingConfig(
        ScalingStrategy strategy, Integer minReplicas, Integer maxReplicas,
        List<ScalingMetric> metrics, ConcurrencyControlConfig concurrencyControl
) { ... }
\end{lstlisting}

L'annidamento merita una difesa, perch\'e sembra contraddire l'affermazione del
\cref{cap:02-stato-dell-arte} secondo cui \nf{} separa la decisione sul numero di
repliche da quella sul limite di concorrenza. La separazione \`e sul piano dei
\emph{meccanismi} --- due moduli distinti, sostituibili indipendentemente, che
non si conoscono --- non su quello della configurazione, dove entrambe le
politiche descrivono come una funzione consuma capacit\`a e conviene che siano
dichiarate insieme. Il modo \code{STATIC\_PER\_POD} rende evidente perch\'e i due
piani non possano essere del tutto disgiunti: calcola il limite come
\emph{repliche pronte $\times$ obiettivo per replica}, e quindi legge una
grandezza che l'autoscaler determina.

\section{Il punto di estensione dei moduli}

L'interfaccia che definisce un modulo del control plane \`e minima:

\begin{lstlisting}[language=Java]
public interface ControlPlaneModule {
    default String name() { return getClass().getSimpleName(); }
    Set<Class<?>> configurationClasses();
}
\end{lstlisting}

Un modulo dichiara soltanto quali classi di configurazione Spring devono essere
aggiunte al contesto. Non dichiara dipendenze, non dichiara un ordine, non
espone un ciclo di vita. Tutto ci\`o che gli serve lo ottiene dal contenitore
attraverso i punti di estensione del nucleo. Il \cref{cap:08-sistema-moduli}
descrive il meccanismo di caricamento e discute perch\'e questa povert\`a
dell'interfaccia sia un pregio e non una lacuna.

\section{Il contratto dell'handler}

Infine, il contratto che il codice utente implementa:

\begin{lstlisting}[language=Java]
public interface FunctionHandler {
    Object handle(InvocationRequest request);
}
\end{lstlisting}

Un solo metodo, un solo argomento, nessuna eccezione dichiarata, nessun contesto
di esecuzione iniettato. La povert\`a \`e voluta: qualunque cosa l'handler
riceva come contesto diventerebbe parte del contratto e vincolerebbe gli SDK
degli altri linguaggi a riprodurla. Ci\`o che l'handler pu\`o restituire \`e
governato da \code{HandlerResponse} e \code{ResponseHeaderPolicy}, discussi nel
\cref{cap:17-runtime-funzione}.

\section{L'API come artefatto versionato}

La descrizione formale dell'API vive in \code{openapi.yaml} nella radice del
repository. \`E mantenuta a mano, e questa scelta ha un costo che il progetto ha
sperimentato: essendo esterna al codice, pu\`o divergere in silenzio. La verifica
automatica esistente accerta che il file sia presente, non che sia allineato.

\begin{damisurare}
Il disallineamento fra \code{openapi.yaml} e il codice \`e stato rilevato e
corretto manualmente in almeno un'occasione. Una verifica che confronti gli
endpoint dichiarati nel file con quelli effettivamente esposti dal contesto
Spring \`e realizzabile con gli strumenti gi\`a presenti nel progetto, ma non \`e
stata implementata. \todomis{copertura effettiva di \code{openapi.yaml} rispetto
alle rotte registrate}
\end{damisurare}
